
\documentclass{article}
\usepackage{colortbl}
\usepackage{makecell}
\usepackage{multirow}
\usepackage{supertabular}

\begin{document}

\newcounter{utterance}

\centering \large Interaction Transcript for game `cladder', experiment `full\_v1.5\_default', episode 2504 with selfrefine{-}Qwen3.5{-}27B.
\vspace{24pt}

{ \footnotesize  \setcounter{utterance}{1}
\setlength{\tabcolsep}{0pt}
\begin{supertabular}{c@{$\;$}|p{.15\linewidth}@{}p{.15\linewidth}p{.15\linewidth}p{.15\linewidth}p{.15\linewidth}p{.15\linewidth}}
    \# & \multicolumn{2}{c}{Player} && \multicolumn{2}{c}{Game Master} \\
    \hline

    \theutterance \stepcounter{utterance}  
    & & & \multicolumn{4}{p{0.6\linewidth}}{
        \cellcolor[rgb]{0.9,0.9,0.9}{
            \makecell[{{p{\linewidth}}}]{
                \texttt{\tiny{[P1$\langle$GM]}}
                \texttt{You are an expert in causal inference. The following question is not a typical commonsense query, but rather a meticulously designed question created by a professor specializing in causal inference, intended to assess the students' mastery of the course content.} \\
\texttt{Imagine a self{-}contained, hypothetical world with only the following conditions, and without any unmentioned factors or causal relationships: Kidney stone size has a direct effect on treatment and recovery. Treatment has a direct effect on recovery. For patients who have small kidney stones and not receiving treatment, the probability of recovery is 88\%. For patients who have small kidney stones and receiving treatment, the probability of recovery is 72\%. For patients who have large kidney stones and not receiving treatment, the probability of recovery is 16\%. For patients who have large kidney stones and receiving treatment, the probability of recovery is 4\%. The overall probability of large kidney stone is 41\%. For patients receiving treatment, would it be less likely to see recovery if the patient had received no treatment?} \\
\texttt{Start your answer with 'Yes' or 'No', followed by additional reasoning or evidence to support your explanation.} \\
            }
        }
    }
     \\ \\

    \theutterance \stepcounter{utterance}  
    & & \multicolumn{4}{p{0.6\linewidth}}{
        \cellcolor[rgb]{0.95,0.95,0.95}{
            \makecell[{{p{\linewidth}}}]{
                \texttt{\tiny{[P1$\leftrightarrow$P]}}
                \texttt{{[}DRAFT{]}} \\
\texttt{No} \\
\\ 
\texttt{The reasoning relies on the distinction between **marginal (observed) association** and **causal effect**, specifically illustrating **Simpson's Paradox** (or the need to adjust for a confounder).} \\
\\ 
\texttt{Here is the step{-}by{-}step causal analysis:} \\
\\ 
\texttt{1.  **Identify the Causal Structure**:} \\
\texttt{    *   **Confounder (\$Z\$)**: Kidney stone size. It affects both the choice of **Treatment (\$T\$)** and the **Recovery (\$Y\$)**.} \\
\texttt{    *   **Treatment (\$T\$)**: Affects Recovery (\$Y\$).} \\
\texttt{    *   **Recovery (\$Y\$)**: The outcome.} \\
\texttt{    *   The question asks for the **Average Causal Effect (ACE)** of treatment: \$P(Y=1 | do(T=1)) {-} P(Y=1 | do(T=0))\$.} \\
\\ 
\texttt{2.  **Analyze the Conditional Probabilities (Stratified by Stone Size)**:} \\
\texttt{    *   **Small Stones (\$Z=S\$)**:} \\
\texttt{        *   \$P(Y=1 | T=1, Z=S) = 0.72\$} \\
\texttt{        *   \$P(Y=1 | T=0, Z=S) = 0.88\$} \\
\texttt{        *   *Effect*: Treatment reduces recovery probability by \$0.72 {-} 0.88 = {-}0.16\$.} \\
\texttt{    *   **Large Stones (\$Z=L\$)**:} \\
\texttt{        *   \$P(Y=1 | T=1, Z=L) = 0.04\$} \\
\texttt{        *   \$P(Y=1 | T=0, Z=L) = 0.16\$} \\
\texttt{        *   *Effect*: Treatment reduces recovery probability by \$0.04 {-} 0.16 = {-}0.12\$.} \\
\\ 
\texttt{3.  **Calculate the Causal Effect**:} \\
\texttt{    Since the effect of treatment is negative in *both* strata (small and large stones), the overall causal effect of treatment must be negative.} \\
\texttt{    *   The probability of recovery if treated (\$do(T=1)\$) is a weighted average of the treated probabilities:} \\
\texttt{        \$\$P(Y=1 | do(T=1)) = P(Y=1 | T=1, S)P(S) + P(Y=1 | T=1, L)P(L)\$\$} \\
\texttt{        \$\$= 0.72 \textbackslash{}times P(S) + 0.04 \textbackslash{}times P(L)\$\$} \\
\texttt{    *   The probability of recovery if not treated (\$do(T=0)\$) is a weighted average of the untreated probabilities:} \\
\texttt{        \$\$P(Y=1 | do(T=0)) = P(Y=1 | T=0, S)P(S) + P(Y=1 | T=0, L)P(L)\$\$} \\
\texttt{        \$\$= 0} \\
            }
        }
    }
    & \\ \\

    \theutterance \stepcounter{utterance}  
    & & \multicolumn{4}{p{0.6\linewidth}}{
        \cellcolor[rgb]{0.95,0.95,0.95}{
            \makecell[{{p{\linewidth}}}]{
                \texttt{\tiny{[P1$\leftrightarrow$P]}}
                \texttt{{[}SELF{-}FEEDBACK{]}} \\
\texttt{FORMAT: problem} \\
\texttt{RULES: ok} \\
\texttt{STRATEGY: ok} \\
\texttt{REASON: The answer draft is incomplete and cuts off mid{-}calculation, failing to provide the full reasoning required by the prompt.} \\
            }
        }
    }
    & \\ \\

    \theutterance \stepcounter{utterance}  
    & \multicolumn{4}{p{0.6\linewidth}}{
        \cellcolor[rgb]{0.561,0.737,0.561}{
            \makecell[{{p{\linewidth}}}]{
                \texttt{\tiny{[P1$\rangle$GM]}}
                \texttt{No} \\
\\ 
\texttt{The reasoning relies on the distinction between **marginal (observed) association** and **causal effect**, specifically illustrating **Simpson's Paradox** (or the need to adjust for a confounder).} \\
\\ 
\texttt{Here is the step{-}by{-}step causal analysis:} \\
\\ 
\texttt{1.  **Identify the Causal Structure**:} \\
\texttt{    *   **Confounder (\$Z\$)**: Kidney stone size. It affects both the choice of **Treatment (\$T\$)** and the **Recovery (\$Y\$)**.} \\
\texttt{    *   **Treatment (\$T\$)**: Affects Recovery (\$Y\$).} \\
\texttt{    *   **Recovery (\$Y\$)**: The outcome.} \\
\texttt{    *   The question asks for the **Average Causal Effect (ACE)** of treatment: \$P(Y=1 | do(T=1)) {-} P(Y=1 | do(T=0))\$.} \\
\\ 
\texttt{2.  **Analyze the Conditional Probabilities (Stratified by Stone Size)**:} \\
\texttt{    *   **Small Stones (\$Z=S\$)**:} \\
\texttt{        *   \$P(Y=1 | T=1, Z=S) = 0.72\$} \\
\texttt{        *   \$P(Y=1 | T=0, Z=S) = 0.88\$} \\
\texttt{        *   *Effect*: Treatment reduces recovery probability by \$0.72 {-} 0.88 = {-}0.16\$.} \\
\texttt{    *   **Large Stones (\$Z=L\$)**:} \\
\texttt{        *   \$P(Y=1 | T=1, Z=L) = 0.04\$} \\
\texttt{        *   \$P(Y=1 | T=0, Z=L) = 0.16\$} \\
\texttt{        *   *Effect*: Treatment reduces recovery probability by \$0.04 {-} 0.16 = {-}0.12\$.} \\
\\ 
\texttt{3.  **Calculate the Causal Effect**:} \\
\texttt{    Since the effect of treatment is negative in *both* strata (small and large stones), the overall causal effect of treatment must be negative, regardless of the distribution of stone sizes.} \\
\texttt{    *   The probability of recovery if treated (\$do(T=1)\$) is a weighted average of the treated probabilities:} \\
\texttt{        \$\$P(Y=1 | do(T=1)) = P(Y=1 | T=1, S)P(S) + P(Y=1 | T=1, L)P(L)\$\$} \\
\texttt{        \$\$= 0.72 \textbackslash{}times P(S) + 0.04 \textbackslash{}times P(L)\$\$} \\
\texttt{    *   The probability of recovery if not treated (\$do(T=0)\$) is a weighted average of the untreated probabilities:} \\
\texttt{        \$\$P(Y=1 | do(T=0)) = P(Y=1 | T=0, S)P(S) + P(Y=1 | T=0, L)P(L} \\
            }
        }
    }
    & & \\ \\

    \theutterance \stepcounter{utterance}  
    & & & \multicolumn{2}{p{0.3\linewidth}}{
        \cellcolor[rgb]{0.9,0.9,0.9}{
            \makecell[{{p{\linewidth}}}]{
                \texttt{\tiny{[GM$|$GM]}}
                \texttt{No} \\
            }
        }
    }
    & & \\ \\

    \theutterance \stepcounter{utterance}  
    & & & \multicolumn{2}{p{0.3\linewidth}}{
        \cellcolor[rgb]{0.9,0.9,0.9}{
            \makecell[{{p{\linewidth}}}]{
                \texttt{\tiny{[GM$|$GM]}}
                \texttt{no} \\
            }
        }
    }
    & & \\ \\

    \theutterance \stepcounter{utterance}  
    & & & \multicolumn{2}{p{0.3\linewidth}}{
        \cellcolor[rgb]{0.9,0.9,0.9}{
            \makecell[{{p{\linewidth}}}]{
                \texttt{\tiny{[GM$|$GM]}}
                \texttt{game\_result = WIN} \\
            }
        }
    }
    & & \\ \\

\end{supertabular}
}

\end{document}
